\RequirePackage{currfile}
\documentclass[12pt]{beamer}
\usepackage[utf8]{inputenc}
\usepackage[spanish]{babel}
\usepackage{standalone}
\usepackage{color}
\usepackage{siunitx}
\usepackage{hyperref}
%\hypersetup{colorlinks,linkcolor=,urlcolor=blue}
%\hypersetup{colorlinks,urlcolor=blue}
\usepackage{xcolor,soul}
\usepackage{etoolbox}
\usepackage{amsmath}
\usepackage{amsthm}
\usepackage{physics}
\usepackage{multicol}
\usepackage{bookmark}
\usepackage{longtable}
\usepackage{listings}
\usepackage{graphicx}
\usepackage{tikz}
\usetikzlibrary{patterns, matrix, backgrounds, decorations,shapes, arrows.meta}
\usepackage[autostyle,spanish=mexican]{csquotes}
\usepackage[os=win]{menukeys}
\usepackage{pifont}
\usepackage{pbox}
\usepackage{caption}
\captionsetup{font=scriptsize,labelfont=scriptsize}
%\usepackage[sfdefault]{roboto}  %% Option 'sfdefault' only if the base font of the document is to be sans serif

%Sección de definición de colores
\definecolor{ao}{rgb}{0.0, 0.5, 0.0}
\definecolor{bisque}{rgb}{1.0, 0.89, 0.77}
\definecolor{amber}{rgb}{1.0, 0.75, 0.0}
\definecolor{armygreen}{rgb}{0.29, 0.33, 0.13}
\definecolor{alizarin}{rgb}{0.82, 0.1, 0.26}
\definecolor{cadetblue}{rgb}{0.37, 0.62, 0.63}
\definecolor{deepblue}{rgb}{0,0,0.5}
\definecolor{brown}{rgb}{0.59, 0.29, 0.0}
\definecolor{OliveGreen}{rgb}{0,0.25,0}


\usefonttheme[onlymath]{serif}
%Sección de definición de nuevos comandos

\newcommand*{\TitleParbox}[1]{\parbox[c]{1.75cm}{\raggedright #1}}%
\newcommand{\python}{\texttt{python}}
\newcommand{\textoazul}[1]{\textcolor{blue}{#1}}
\newcommand{\azulfuerte}[1]{\textcolor{blue}{\textbf{#1}}}
\newcommand{\funcionazul}[1]{\textcolor{blue}{\textbf{\texttt{#1}}}}
\newcommand{\ptilde}[1]{\ensuremath{{#1}^{\prime}}}
\newcommand{\stilde}[1]{\ensuremath{{#1}^{\prime \prime}}}
\newcommand{\ttilde}[1]{\ensuremath{{#1}^{\prime \prime \prime}}}
\newcommand{\ntilde}[2]{\ensuremath{{#1}^{(#2)}}}
\renewcommand{\arraystretch}{1.5}

\newcounter{saveenumi}
\newcommand{\seti}{\setcounter{saveenumi}{\value{enumi}}}
\newcommand{\conti}{\setcounter{enumi}{\value{saveenumi}}}
\renewcommand{\rmdefault}{cmr}% cmr = Computer Modern Roman

\linespread{1.5}

\usefonttheme{professionalfonts}
%\usefonttheme{serif}
\DeclareGraphicsExtensions{.pdf,.png,.jpg}


%Sección para el tema de beamer, con el theme, usercolortheme y sección de footers
\mode<presentation>
{
  \usetheme{Warsaw}
  
  %\useoutertheme{infolines}
  \useoutertheme{default}
  \usecolortheme{spruce}
  \setbeamercovered{invisible}
  % or whatever (possibly just delete it)
  \setbeamertemplate{section in toc}[sections numbered]
  \setbeamertemplate{subsection in toc}[subsections numbered]
  \setbeamertemplate{subsection in toc}{\leavevmode\leftskip=3.2em\rlap{\hskip-2em\inserttocsectionnumber.\inserttocsubsectionnumber}\inserttocsubsection\par}
  \setbeamercolor{section in toc}{fg=blue}
  \setbeamercolor{subsection in toc}{fg=blue}
  \setbeamercolor{frametitle}{fg=blue}
  \setbeamertemplate{caption}[numbered]

  \setbeamertemplate{footline}
  \beamertemplatenavigationsymbolsempty
  \setbeamertemplate{headline}{}
}

\makeatletter
\setbeamercolor{section in foot}{bg=gray!30, fg=black!90!orange}
\setbeamercolor{subsection in foot}{bg=blue!30!yellow, fg=red}
\setbeamertemplate{footline}
{
  \leavevmode%
  \hbox{%
  \begin{beamercolorbox}[wd=.333333\paperwidth,ht=2.25ex,dp=1ex,center]{section in foot}%
    \usebeamerfont{section in foot} \insertsection
  \end{beamercolorbox}}%
  \begin{beamercolorbox}[wd=.333333\paperwidth,ht=2.25ex,dp=1ex,center]{subsection in foot}%
    \usebeamerfont{subsection in foot}  \insertsubsection
  \end{beamercolorbox}%
  \begin{beamercolorbox}[wd=.333333\paperwidth,ht=2.25ex,dp=1ex,right]{date in head/foot}%
    \usebeamerfont{date in head/foot} \insertshortdate{} \hspace*{2em}
    \insertframenumber{} / \inserttotalframenumber \hspace*{2ex} 
  \end{beamercolorbox}}%
  \vskip0pt%
\makeatother  

\makeatletter
\patchcmd{\beamer@sectionintoc}
  {\vfill}
  {\vskip\itemsep}
  {}
  {}
\makeatother


\title{\large{Ejercicios}}
\subtitle{Funciones Gamma y Beta}
\author{M. en C. Gustavo Contreras Mayén}
\date{\today}
\institute{Facultad de Ciencias - UNAM}
\titlegraphic{\includegraphics[width=1.75cm]{../Imagenes/escudo-facultad-ciencias}\hspace*{4.75cm}~%
   \includegraphics[width=1.75cm]{../Imagenes/escudo-unam}
}
\setbeamertemplate{navigation symbols}{}
\begin{document}
\maketitle
\fontsize{14}{14}\selectfont
\spanishdecimal{.}
\section*{Contenido}
\frame[allowframebreaks]{\tableofcontents[currentsection, hideallsubsections]}
\section{Función Gamma}
\frame{\tableofcontents[currentsection, hideothersubsections]}
\subsection{Identidad}
\begin{frame}
\frametitle{Propiedad función $\Gamma(x)$}
\textbf{Demuestra} la siguiente identidad:
\begin{align*}
\Gamma (x + 1) = x \, \Gamma (x)
\end{align*}
\end{frame}
\begin{frame}
\frametitle{Solución}
Antes de demostrar la identidad directamente de la integral Gamma para todo $x$ positivo, veamos que esta identidad se usa para definir la función Gamma primero para $-1 < x < 0$ escribiéndola en la forma 
\begin{align*}
\Gamma(x) = \dfrac{\Gamma (x + 1)}{x}
\end{align*}
\end{frame}
\begin{frame}
\frametitle{Solución}
Tendremos entonces que la expresión es válida para el intervalo $-2 < x < -1$, y así sucesivamente para todos los valores no enteros negativos de $x$.
\\
\bigskip
\pause
Por lo que nos queda demostrar que se cumple
\begin{align*}
\Gamma (x + 1) =  x \, \Gamma (x)
\end{align*}
para todo $x$ positivo.
\end{frame}
\begin{frame}
\frametitle{Demostración}
Haciendo que $x$ sea cualquier número positivo, escribimos la integral Gamma para el argumento $x + 1$:
\begin{eqnarray*}
\Gamma (x + 1) &= \displaystyle \int_{0}^{\infty} t^{(x+1)-1} \, e^{-t} \dd{t} = \\[0.5em] \pause
&= t^{x} \, e^{-t} \dd{t}
\end{eqnarray*}
\end{frame}
\begin{frame}
\frametitle{Demostración}
Resolvemos usando la integración por partes, haciendo:
\begin{align*}
u = t^{x} \hspace{1cm} \dd{v} = e^{-t} \dd{t}
\end{align*}
entonces
\pause
\begin{align*}
\Gamma (x + 1) = t^{x} \, \left( - e^{-t}\right)\eval_{0}^{\infty} - \int_{0}^{\infty} \left( -e^{-t}\right) \, x \, t^{x-1} \dd{t}
\end{align*}
\end{frame}
\begin{frame}
\frametitle{Demostración}
Por lo que:
\begin{align*}
\Gamma (x + 1) = - \lim_{t \to \infty} \dfrac{t^{x}}{e^t} + \dfrac{0^{x}}{e^{0}} + x \int_{0}^{\infty} t^{x-1} \, e^{-t} \dd{t}
\end{align*}
\pause
Veamos ahora lo que pasa con este resultado:
\end{frame}
\begin{frame}
\frametitle{Demostración}
El límite en el primer sumando de la expresión de la derecha sabemos que se anula, a partir de la indeterminación $\infty / \infty$, ya que usamos la regla de L'Hopital.
\begin{align*}
\Gamma (x + 1) = \cancelto{0}{- \lim_{t \to \infty} \dfrac{t^{x}}{e^t}} + \dfrac{0^{x}}{e^{0}} + x \int_{0}^{\infty} t^{x-1} \, e^{-t} \dd{t}
\end{align*}
\end{frame}
\begin{frame}
\frametitle{Demostración}
El segundo término de la derecha también se anula
\begin{align*}
\Gamma (x + 1) = \cancelto{0}{- \lim_{t \to \infty} \dfrac{t^{x}}{e^t}} + \cancelto{0}{\dfrac{0^{x}}{e^{0}}} + x \int_{0}^{\infty} t^{x-1} \, e^{-t} \dd{t}
\end{align*}
\pause
Por lo que nos queda el tercero, pero vemos que la integral que queda, es la definición de la función Gamma, así:
\end{frame}
\begin{frame}
\frametitle{Demostración}
Tenemos que
\begin{align*}
\Gamma (x + 1) = x \, \Gamma (x) \qed
\end{align*}
\end{frame}
\end{document}